% !TeX program = xelatex
% UTF-8; standalone XeLaTeX/XePersian transcription.
% Source: Google Drive / Notebook / MSPS. Source statements are not clinically revised.
% images/pic01.jpg = 20261018_193208391.jpg (1TSvRcnER1nL1pQUZ3n5dtxgsfkpxvQlX)
% images/pic02.jpg = 20261018_193213092.jpg (1zm3hfiWdJPZOzoTeZDhUTGjFbmBL0Ely)
% Printed writing-guide labels on both images: 1, 2, 3, 4, 5, 6, 7, 8, 9.
% Handwritten footer numbers on the two images: ۱ and ۳, respectively.
\documentclass[11pt,a4paper]{article}
\usepackage[margin=18mm]{geometry}
\usepackage{amsmath,amssymb}
\usepackage{enumitem}
\usepackage{xcolor}
\usepackage{xepersian}
\settextfont{Amiri}
\setlatintextfont{DejaVu Serif}
\setlength{\parindent}{0pt}
\setlength{\parskip}{3pt}
\setlength{\emergencystretch}{3em}
\setlist[itemize]{itemsep=2pt,topsep=3pt,parsep=0pt}
\newcommand{\en}[1]{\lr{#1}}
\newcommand{\hl}[1]{\colorbox{yellow!45}{\strut #1}}
\newcommand{\arr}{\lr{\(\leftarrow\)}}
\newcommand{\sourcehead}{%
  \noindent\en{Subject:} ۹۲ ـ سندرم‌های درد اسکلتی عضلانی\par
  \noindent\en{Year: \ldots\qquad Month: \ldots\qquad Date: \ldots}\par
  \smallskip\hrule\medskip}
\begin{document}
\sourcehead

\subsection*{درد رشد}
\begin{itemize}
\item در ۴۰\% کودکان سن مدرسه \arr{} \hl{۳–۱۲ سال}؛ در \hl{پسرها} شایع‌تر است.
\item نوع درد: عمقی، کرامپی، در ساق پا؛ \hl{در ناحیهٔ مفاصل نیست}.
\item \hl{عصر و شب بیشتر رخ می‌دهد}؛ می‌تواند کودک را از خواب بیدار کند.
\item \hl{با فعالیت بدنی تشدید می‌شود}.
\item آزمایش‌ها و تصویربرداری‌ها \en{N/L} هستند.
\item در برخورد با کودک با این شکایت \arr{} \hl{\en{R/O} لوکمی} \arr{} \hl{\en{CBC}}.
\item درمان \arr{} استامینوفن و \en{NSAID} در صورت نیاز.
\end{itemize}

\medskip\hrule\medskip
\subsection*{سندرم‌های درد مایوفاشیال و فیبرومیالژی}
\begin{itemize}
\item درد اسکلتی عضلانی منتشر، چندین نقطهٔ تندر، خستگی، ضعف، اختلال خواب.
\item در کودکان \hl{بالای ۱۲ سال}؛ \hl{\en{\(M<F\)}}.
\item ممکن است به دنبال تروما یا عفونت باشد \arr{} \hl{آنفولانزا و \en{EBV}}.
\item درد در هر ساعتی از روز می‌تواند باشد.
\item نباید بی‌دلیل \en{ANA} درخواست کنیم \arr{} در ۲\% کودکان مثبت است.
\item تشخیص \arr{} نقاط تندر متعدد در نبود بیماری؛ درد منتشر حداقل به مدت ۳ ماه.
\item درمان دارویی:
  \begin{itemize}
  \item آمی‌تریپتیلین \arr{} برای تنظیم خواب.
  \item گاباپنتین \arr{} کاهش حساسیت به درد.
  \end{itemize}
\end{itemize}
\vfill
\noindent ۱

\clearpage
\sourcehead
\subsection*{هایپرموبیلیتی خوش‌خیم}
\begin{itemize}
\item علائم \arr{} درد اسکلتی عضلانی ثانویه به دامنهٔ حرکت بالای مفاصل؛ درد می‌تواند در روز یا شب باشد.
\item کودکان \hl{۳–۱۰ سال}؛ \hl{\en{\(M<F\)}}.
\item بیماری‌هایی که منجر به هایپرموبیلیتی می‌شوند:
  سندرم مارفان، هوموسیستین‌اوری، سندرم استیکلر، سندرم اهلرز–دانلوس،
  استئوژنزيس ایمپرفکتا، سندرم ویلیامز، تریزومی ۲۱.
\item می‌تواند جنرالیزه و یا محدود به گروه خاصی از مفاصل باشد.
\item \en{genu recurvatum} \arr{} افزایش \en{ROM} مفاصل \en{MCP}، مچ، آرنج، زانو.
\item \hl{\en{pes planus}} \arr{} \hl{\en{flat foot}}.
\item اگر هایپرموبیلیتی با \hl{الاستیسیتهٔ بالای پوست}، \hl{ایزی برویزینگ}، \hl{\en{MVP}} همراه باشد
  \arr{} شک به \hl{اهلرز–دانلوس، مارفان}.
\end{itemize}
\vfill
\noindent ۳
\end{document}
